\documentclass{article}
\usepackage{../style/latex/dhruv_theory_style}

% -------------------------------------------------------
% Document metadata
% -------------------------------------------------------
\title{\textbf{DSC 190/291 Assignment 1}}
\author{Dhruv Patel}
\date{\today}

% =======================================================
\begin{document}
% =======================================================

\maketitle
\tableofcontents
\newpage

% =======================================================
\section{Part A: Repository Setup}
% =======================================================
Public GitHub repository: \href{https://github.com/PDhruv09/dsc190_learning_theory}{github.com/PDhruv09/dsc190_learning_theory}
% =======================================================
\section{Part B: Theory Problem}
% =======================================================
\subsection{Problem Overview}

This problem starts from the threshold story in lecture and then pushes it one
step further.

Throughout, let $\mathcal{X} = [0,1]$, $\mathcal{Y} = \{-1,+1\}$, and let
\[
h_\theta(x) = \sign(x-\theta),
\]
where we use the convention $\sign(z)=+1$ for $z \ge 0$ and $\sign(z)=-1$ for
$z<0$.

We will study sequences with one extra structural assumption beyond what was
explicitly analyzed in lecture.
\newline

\textbf{$\Delta$-Separated Threshold-Realizable Sequences.}
We say a sequence $((x_t,y_t))_{t=1}^T$ is \emph{$\Delta$-separated
threshold-realizable} if there exist $\theta^* \in [0,1]$ and $\Delta > 0$
such that
\[
y_t = h_{\theta^*}(x_t)
\qquad \text{and} \qquad
|x_t-\theta^*| \ge \Delta
\qquad \text{for every round } t.
\]

\qquad
\qquad

\subsection{From Continuous Thresholds to a Finite Class}
% =======================================================

\textbf{Question.}
Design a finite grid $G \subseteq [0,1]$ whose size depends only on $\Delta$,
and consider the associated finite threshold class
\[
H_G = \{h_\theta : \theta \in G\}.
\]
Prove that for every $\Delta$-separated threshold-realizable sequence, there
exists some $\theta \in G$ such that
\[
h_\theta(x_t)=y_t
\qquad \text{for all } t=1,\dots,T.
\]
Then use the Halving theorem from lecture to derive a mistake bound of order
\[
O\!\left(\log\left(\frac{1}{\Delta}\right)\right)
\]
for an explicit online learner.

Your answer should clearly state:
\begin{itemize}[leftmargin=2.5em]
    \item your choice of grid $G$,
    \item the size of $G$ as a function of $\Delta$,
    \item the resulting mistake bound.
\end{itemize}

\begin{solutionbox}

\subsubsection{Construction of the Grid}

We choose a grid with spacing $\Delta/2$:
\[
G
=
\left\{
k\cdot \frac{\Delta}{2}
:
k=0,1,2,\dots,\left\lfloor \frac{2}{\Delta}\right\rfloor
\right\}
\cap [0,1].
\]

Its size satisfies
\[
|G|
\le
\left\lfloor \frac{2}{\Delta}\right\rfloor + 1
=
O\!\left(\frac{1}{\Delta}\right).
\]

The reason for spacing $\Delta/2$ is that for every $\theta^* \in [0,1]$,
there exists some grid point $\theta \in G$ such that
\[
|\theta-\theta^*| \le \frac{\Delta}{2}.
\]

\subsubsection{Proof That a Grid Threshold Is Consistent}

\begin{claim}
\leavevmode\\
    For every $\Delta$-separated threshold-realizable sequence
    $(x_t,y_t)_{t=1}^T$, there exists $\theta \in G$ such that
    \[
    h_\theta(x_t)=y_t
    \qquad \text{for all } t=1,\dots,T.
    \]
\end{claim}
\begin{proof}
\leavevmode\\
Because the sequence is $\Delta$-separated threshold-realizable, there exists
$\theta^* \in [0,1]$ such that
\[
y_t = h_{\theta^*}(x_t)=\sign(x_t-\theta^*)
\qquad \text{for all } t,
\]
and
\[
|x_t-\theta^*| \ge \Delta
\qquad \text{for all } t.
\]

By construction of the grid, there exists some $\theta \in G$ with
\[
|\theta-\theta^*| \le \frac{\Delta}{2}.
\]

We now show that $h_\theta(x_t)=h_{\theta^*}(x_t)$ for every $t$.

\textbf{Case 1:} $y_t=+1$.

Then $x_t \ge \theta^*$, and by separation,
\[
x_t-\theta^* \ge \Delta.
\]
Using $|\theta-\theta^*| \le \Delta/2$,
\[
x_t-\theta
=
(x_t-\theta^*)-(\theta-\theta^*)
\ge
\Delta-\frac{\Delta}{2}
=
\frac{\Delta}{2}
>0.
\]
So $x_t \ge \theta$, hence
\[
h_\theta(x_t)=+1=y_t.
\]

\textbf{Case 2:} $y_t=-1$.

Then $x_t<\theta^*$, and by separation,
\[
\theta^*-x_t \ge \Delta.
\]
Again using $|\theta-\theta^*| \le \Delta/2$,
\[
\theta-x_t
=
(\theta^*-x_t)-(\theta^*-\theta)
\ge
\Delta-\frac{\Delta}{2}
=
\frac{\Delta}{2}
>0.
\]
So $x_t<\theta$, hence
\[
h_\theta(x_t)=-1=y_t.
\]

Thus $h_\theta(x_t)=y_t$ for every $t$.
\end{proof}

\subsubsection{Halving Algorithm and Mistake Bound}

Consider the finite class
\[
H_G=\{h_\theta:\theta\in G\}.
\]
Run the Halving algorithm on this class.

Since the claim above shows that there exists a hypothesis in $H_G$ consistent
with the entire sequence, the Halving theorem applies. Therefore the number of
mistakes is at most
\[
M \le \log_2 |H_G| = \log_2 |G|.
\]
Using the bound on the size of the grid,
\[
M
\le
\log_2\!\left(\left\lfloor \frac{2}{\Delta}\right\rfloor+1\right)
=
O\!\left(\log\left(\frac{1}{\Delta}\right)\right).
\]

\subsection*{Final Answer for Part 1}

\begin{itemize}[leftmargin=2.5em]
    \item \textbf{Choice of grid:}
    \[
    G=
    \left\{
    k\cdot \frac{\Delta}{2}
    :
    k=0,1,2,\dots,\left\lfloor \frac{2}{\Delta}\right\rfloor
    \right\}
    \cap [0,1].
    \]

    \item \textbf{Size of the grid:}
    \[
    |G|
    \le
    \left\lfloor \frac{2}{\Delta}\right\rfloor+1
    =
    O\!\left(\frac{1}{\Delta}\right).
    \]

    \item \textbf{Mistake bound:}
    \[
    M
    \le
    \log_2\!\left(\left\lfloor \frac{2}{\Delta}\right\rfloor+1\right)
    =
    O\!\left(\log\left(\frac{1}{\Delta}\right)\right).
    \]
\end{itemize}

\end{solutionbox}

\subsection{A Positive Margin from Separation}
\textbf{Question.}
\newline
View thresholds as linear predictors by choosing a feature map
\[
\phi:[0,1]\to \mathbb{R}^d
\]
and a unit vector $u^*$ depending on $\theta^*$.

Prove that every $\Delta$-separated threshold-realizable sequence is linearly
separable with margin at least $c\Delta$ for some absolute constant $c>0$
under your representation, while
\[
\|\phi(x)\|\le R
\qquad \text{for all } x\in [0,1]
\]
for some absolute constant $R$.

Then use the Perceptron theorem from lecture to derive an explicit mistake
bound of order
\[
O\!\left(\frac{1}{\Delta^2}\right).
\]

Your answer should clearly specify your chosen feature map, the margin lower
bound, the norm bound, and the final mistake bound.

\begin{solutionbox}

\subsubsection{Feature Map and Linear Representation}

Chosen
\[
\phi(x)=(x,1)\in \R^2.
\]
Let
\[
w^*=(1,-\theta^*).
\]
Then
\[
\langle w^*,\phi(x)\rangle
=
1\cdot x + (-\theta^*)\cdot 1
=
x-\theta^*.
\]
Therefore
\[
\sign(\langle w^*,\phi(x)\rangle)
=
\sign(x-\theta^*)
=
h_{\theta^*}(x).
\]

So this representation realizes threshold classification exactly.

Now normalize:
\[
u^*=\frac{w^*}{\|w^*\|}.
\]
Since $\theta^*\in[0,1]$,
\[
\|w^*\|
=
\sqrt{1+(\theta^*)^2}
\le
\sqrt{2}.
\]

\subsubsection{Margin Lower Bound}

For every round $t$,
\[
y_t\langle u^*,\phi(x_t)\rangle
=
\frac{y_t\langle w^*,\phi(x_t)\rangle}{\|w^*\|}
=
\frac{y_t(x_t-\theta^*)}{\|w^*\|}.
\]
Because
\[
y_t=\sign(x_t-\theta^*),
\]
we get
\[
y_t(x_t-\theta^*) = |x_t-\theta^*|.
\]
By the $\Delta$-separation assumption,
\[
|x_t-\theta^*| \ge \Delta.
\]
Also $\|w^*\|\le \sqrt{2}$. Hence
\[
y_t\langle u^*,\phi(x_t)\rangle
\ge
\frac{\Delta}{\sqrt{2}}.
\]

Thus the sequence is linearly separable with margin at least
\[
\gamma \ge \frac{\Delta}{\sqrt{2}}.
\]
So we can take
\[
c=\frac{1}{\sqrt{2}}.
\]

\subsubsection{Norm Bound}

For any $x\in[0,1]$,
\[
\|\phi(x)\|^2 = x^2+1 \le 2.
\]
Therefore
\[
\|\phi(x)\| \le \sqrt{2}.
\]
So we may take
\[
R=\sqrt{2}.
\]

\subsubsection{Perceptron Mistake Bound}

By the Perceptron theorem, the number of mistakes satisfies
\[
M \le \left(\frac{R}{\gamma}\right)^2.
\]
Substituting the bounds above,
\[
M
\le
\left(\frac{\sqrt{2}}{\Delta/\sqrt{2}}\right)^2
=
\left(\frac{2}{\Delta}\right)^2
=
\frac{4}{\Delta^2}.
\]
Hence
\[
M = O\!\left(\frac{1}{\Delta^2}\right).
\]

\subsubsection*{Final Answer for Part 2}

\begin{itemize}[leftmargin=2.5em]
    \item \textbf{Feature map:}
    \[
    \phi(x)=(x,1)\in\R^2.
    \]

    \item \textbf{Margin lower bound:}
    \[
    \gamma \ge \frac{\Delta}{\sqrt{2}}.
    \]

    \item \textbf{Norm bound:}
    \[
    \|\phi(x)\|\le \sqrt{2}
    \qquad \text{for all } x\in[0,1].
    \]

    \item \textbf{Mistake bound:}
    \[
    M\le \frac{4}{\Delta^2}
    =
    O\!\left(\frac{1}{\Delta^2}\right).
    \]
\end{itemize}

\end{solutionbox}

\subsection{Comparison and Interpretation}
\textbf{Question.}
\newline
Explain why the continuous-threshold impossibility phenomenon from lecture does
not contradict Part 1.

Then compare the two mistake bounds you obtained in Parts 1 and 2. Why do they
scale differently with $\Delta$? What is each argument measuring about the
problem?

\begin{solutionbox}

The continuous-threshold impossibility result from lecture applies to the full
class
\[
\mathcal{H}=\{h_\theta:\theta\in[0,1]\}
\]
with no structural restriction on the sequence. In that unrestricted setting,
an adversary can place examples arbitrarily close to the true threshold, which
prevents an online learner from obtaining a finite worst-case mistake bound.

Part 1 does not contradict this because it assumes something extra:
$\Delta$-separation. Once every example must remain at least $\Delta$ away from
the true threshold $\theta^*$, thresholds within distance $\Delta/2$ of
$\theta^*$ become indistinguishable on the observed data. That is why the
continuous class can be replaced by a finite grid. So the impossibility result
and the positive result in Part 1 refer to different settings.

Now compare the two bounds:
\[
\text{Part 1: } O\!\left(\log\left(\frac{1}{\Delta}\right)\right),
\qquad
\text{Part 2: } O\!\left(\frac{1}{\Delta^2}\right).
\]

They scale differently because they come from different kinds of arguments.

The Part 1 argument is \emph{combinatorial}. It counts how many effectively
different hypotheses remain after using the separation assumption. Since the
effective finite class has size
\[
O\!\left(\frac{1}{\Delta}\right),
\]
the Halving algorithm gives a logarithmic mistake bound.

The Part 2 argument is \emph{geometric}. It uses linear separability and the
margin bound. The Perceptron theorem depends on
\[
\left(\frac{R}{\gamma}\right)^2.
\]
Because the margin is proportional to $\Delta$, the mistake bound scales like
$1/\Delta^2$.

So the Halving argument measures how many candidate thresholds need to be
distinguished, while the Perceptron argument measures how hard the problem is
geometrically in terms of margin.

\end{solutionbox}

\subsection{Optional Strengthening: Audit an AI Proof}
\textbf{Question.}
An AI assistant gives the following argument:

\hangindent=2em
\hangafter=0
Because every example is at least $\Delta$ away from the true threshold, continuous thresholds effectively form a class of size $O(1/\Delta)$. Therefore any online learner, including Perceptron, must make at most $O(\log(1/\Delta))$ mistakes on every $\Delta$-separated threshold-realizable sequence.

Identify at least two mathematical problems with this argument. Then write a
corrected statement that is true and prove it.

\begin{solutionbox}

There are at least two mathematical problems with the AI's argument.

\textbf{Problem 1.}
It incorrectly concludes that because the effective class size is
$O(1/\Delta)$, \emph{any} online learner must make only
$O(\log(1/\Delta))$ mistakes. That is not true. A class-size argument gives a
mistake bound only for specific algorithms such as the Halving algorithm. It is
not a universal guarantee for every possible learner.

\textbf{Problem 2.}
It incorrectly applies the same conclusion to Perceptron. The Perceptron
theorem does not use class size. It uses margin and feature norm. In this
problem, the Perceptron bound is
\[
O\!\left(\frac{1}{\Delta^2}\right),
\]
not
\[
O\!\left(\log\left(\frac{1}{\Delta}\right)\right).
\]

\textbf{Corrected statement.}
For every $\Delta>0$, if we run the Halving algorithm on the finite threshold
class
\[
H_G=\{h_\theta:\theta\in G\},
\]
where
\[
G=
\left\{
k\cdot \frac{\Delta}{2}
:
k=0,1,2,\dots,\left\lfloor \frac{2}{\Delta}\right\rfloor
\right\}
\cap [0,1],
\]
then on every $\Delta$-separated threshold-realizable sequence, the number of
mistakes is at most
\[
\log_2\!\left(\left\lfloor \frac{2}{\Delta}\right\rfloor+1\right)
=
O\!\left(\log\left(\frac{1}{\Delta}\right)\right).
\]

\textbf{Proof.}
From Part 1, there exists some $\theta\in G$ such that
\[
h_\theta(x_t)=y_t
\qquad \text{for all } t=1,\dots,T.
\]
So the class $H_G$ contains a hypothesis consistent with the entire sequence.

The Halving algorithm starts with all hypotheses in $H_G$. Whenever it makes a
mistake, at least half of the currently surviving hypotheses are eliminated.
Because a consistent hypothesis is never removed, the version space always
contains at least one hypothesis. Therefore after $M$ mistakes,
\[
1 \le \frac{|G|}{2^M}.
\]
This implies
\[
M \le \log_2 |G|.
\]
Using
\[
|G| \le \left\lfloor \frac{2}{\Delta}\right\rfloor+1,
\]
we obtain
\[
M
\le
\log_2\!\left(\left\lfloor \frac{2}{\Delta}\right\rfloor+1\right)
=
O\!\left(\log\left(\frac{1}{\Delta}\right)\right).
\]

This corrected statement is true for the Halving algorithm specifically. It
does not imply the same bound for arbitrary online learners or for Perceptron.

\end{solutionbox}

% =======================================================
\section{Part C: Perceptron Implementation and Experiments}
% =======================================================
For this part, I implemented the Perceptron algorithm in Python and designed synthetic experiments to test the mistake bound
\[
M \le \frac{R^2}{\gamma^2}.
\]
The goal was to verify experimentally how the number of mistakes depends on the margin \(\gamma\), the norm bound \(R\), and the dimension \(d\).

To generate data where the assumptions of the theorem are known in advance, I fixed a true separator \(w^* = e_1\), the first standard basis vector. I then generated labels \(y_i \in \{-1,+1\}\), chose values \(a_i \in [\gamma, R]\), and set the first coordinate to \(x_i[0] = y_i a_i\). This guarantees that
\[
y_i (w^* \cdot x_i) = y_i x_i[0] = a_i \geq \gamma.
\]
To ensure that every example satisfies \(\|x_i\| = R\), I filled the remaining coordinates with a random direction scaled so that the full vector lies exactly on the sphere of radius \(R\). This construction gave direct control over both the margin and norm assumptions in the theorem.

I also verified the correctness of the implementation with several sanity checks. I checked that all generated examples had norm exactly \(R\), that every example satisfied the margin lower bound, that the Perceptron update rule matched the lecture definition, that the algorithm converged on separable data, and that the observed mistake count never exceeded the theoretical bound in repeated random trials.

\subsection{Experiment 1: Effect of Margin \(\gamma\)}

In the first experiment, I fixed \(R=5\), \(d=2\), and \(n=500\), and varied \(\gamma\) across several values. I measured the average number of Perceptron mistakes over multiple random trials.

\begin{figure}[h!]
    \centering
    \includegraphics[width=\textwidth]{{hw1/plots/exp1_vary_gamma.png}
    \caption{Mistake count versus margin.}
\end{figure}

The results show that as \(\gamma\) becomes smaller, the number of mistakes increases substantially. This is consistent with the theorem, since the bound scales like \(1/\gamma^2\). However, the empirical curve is not perfectly linear as a function of \(1/\gamma^2\). Instead, the experiment supports the qualitative conclusion that smaller margin leads to more mistakes, with growth broadly consistent with the theoretical prediction.

The same figure also compares the actual mistake count to the bound \(R^2/\gamma^2\). The bound is clearly loose in practice: the theoretical value is often much larger than the observed number of mistakes. This makes sense because the theorem gives a worst-case guarantee over all separable sequences, while these experiments use randomly generated data that is often easier than the worst case.

\subsection{Experiment 2: Effect of Dimension \(d\)}

In the second experiment, I fixed \(\gamma=0.5\), \(R=5\), and \(n=500\), and varied the dimension \(d\) over values from 2 up to 500.

\begin{figure}[h!]
    \centering
    \includegraphics[width=0.8\textwidth]{{hw1/plots/exp2_vary_dimension.png}
    \caption{Mistake count versus dimension.}
\end{figure}

The Perceptron mistake bound does not depend on the ambient dimension, only on \(R\) and \(\gamma\). The experimental results are consistent with that theory: the mistake counts stay in the same general range as dimension increases, and there is no clear monotone growth with \(d\). Although the curve is not perfectly flat, the experiment supports the claim that dimension itself is not the main driver of the mistake count once the geometric quantities \(R\) and \(\gamma\) are fixed.

\subsection{Experiment 3: Near-Zero Margin}

In the third experiment, I pushed the margin toward zero by testing smaller and smaller values of \(\gamma\), again with \(R=5\), \(d=2\), and \(n=500\).

\begin{figure}[h!]
    \centering
    \includegraphics[width=\textwidth]{{hw1/plots/exp3_near_zero_margin.png}
    \caption{Mistake count as the margin approaches zero.}
\end{figure}

As \(\gamma \to 0\), the number of mistakes grows rapidly. This matches the theoretical bound, which diverges like \(1/\gamma^2\). The log-scale view makes the trend easier to see, especially because the theoretical bound becomes extremely large compared with the actual mistakes. Empirically, the growth is broadly compatible with the theory, even though it is not an exact match at every point.

This experiment also connects to the threshold counterexample discussed in lecture. When the margin becomes extremely small, examples lie closer and closer to the decision boundary, making the learning problem much harder. In the limiting case where the data is not strictly separable, the Perceptron no longer has a finite mistake guarantee and may fail to converge. So the blow-up in mistakes as \(\gamma \to 0^+\) can be viewed as a warning sign of the non-separable regime.

\subsection{Summary of Findings}

Overall, the experiments support the main ideas of the Perceptron mistake bound. First, smaller margin leads to more mistakes, in line with the \(1/\gamma^2\) scaling from theory. Second, the theoretical bound is useful but loose: it correctly predicts the trend and guarantees convergence, but it overestimates the actual number of mistakes on these random datasets. Third, the experiments are consistent with the fact that the mistake bound is dimension-free. Finally, as the margin approaches zero, the mistake count increases sharply, which connects naturally to the failure of the Perceptron on non-separable data.
% =======================================================
\section{Part D: AI Usage Report}
% =======================================================
For this assignment, I used AI in several parts of the workflow. I used it to help set up the repository structure, write and organize the Perceptron implementation, design the experimental procedure for testing how the mistake count depends on the margin and dimension, and improve the written explanations of the results. I also used AI to debug presentation issues in the plots and to revise the discussion so that the conclusions matched the actual experimental outputs. In addition, I used AI while preparing the theory writeup and organizing the files required for submission.

One concrete example of an AI suggestion I accepted was the idea for generating synthetic data with a known margin and norm bound. The construction fixed a unit separator vector and chose the first coordinate so that each point satisfied the margin condition, while the remaining coordinates were scaled so that every example had norm exactly \(R\). I accepted this suggestion because it gave direct control over both assumptions in the Perceptron mistake bound, which made the experiments easier to interpret and justify mathematically.

One example of an AI suggestion I rejected or modified was the original interpretation of the experimental plots. The first draft claimed that the relationship between mistakes and \(1/\gamma^2\) was clearly linear and that the mistake count was essentially flat across dimensions. After reviewing the actual plots and outputs more carefully, I realized those statements were too strong. I changed the discussion to say that the experiments were broadly consistent with the theoretical trend, but not a perfect confirmation of exact linearity, and that the dimension experiment showed no clear monotone growth rather than a perfectly flat curve. I also changed the plots themselves so that the y-axis scale made the empirical results visible instead of being dominated by the theoretical bound.

I verified the code and results in several ways. First, I checked the implementation logic directly against the Perceptron update rule from lecture. Then I ran sanity checks on the data generator to confirm that the norm of each example was exactly \(R\) and that every point satisfied the required margin condition. I also tested that the Perceptron converged on separable data, that the learned classifier correctly classified the generated training examples, and that the observed number of mistakes never exceeded the theoretical bound in the random trials I ran. Finally, I compared the written conclusions against the actual plots and printed outputs to make sure the discussion accurately described the experimental evidence rather than overstating it.
\end{document}